%%%
    % C++ Chapter biblio starts
    %%%
    %% %% 1
    %% %%
    \bibitem{fluentc++-mostvexparse}Jonathan Boccara - \emph{The Most Vexing Parse: How to Spot It and Fix It Quickly}
    \newline
    Retrieved from: \href{https://www.fluentcpp.com/2018/01/30/most-vexing-parse/}{Fluent C++}
    
    %% %%
    \bibitem{wiki-mosevexparse}\emph{Most Vexing Parse}
    \newline
    Retrieved from: \href{https://en.wikipedia.org/wiki/Most_vexing_parse}{}
    
    %% %%
    \emph{\textbf{C++ chapter bibliography}} 
    \bibitem{fernandes-ruleofzero} R. Martinho Fernandes (2012) - \emph{Rule of Zero}
    \newline
    Retrieved from: \href{https://isocpp.org/blog/2012/11/rule-of-zero}{ISOC++ Blog}
    %%%

    %% %%
    \bibitem{meyers-concern-ruleof0} Scott Meyers (2014) - \emph{A Concern about the Rule of Zero}
    \newline
    Retrieved from: \href{http://scottmeyers.blogspot.fr/2014/03/a-concern-about-rule-of-zero.html}{The View from Aristeia - Scoot Meyers' Activities and Interests}
    %%%
    
    %% %%
    \bibitem{cppcoreguidelines-C67} \emph{A polymorphic class should suppress public copy/move}
    \newline
    Retreived from: \href{https://github.com/isocpp/CppCoreGuidelines/blob/master/CppCoreGuidelines.md#c67-a-polymorphic-class-should-suppress-public-copymove}{Standard Cpp Foundation Github}
    
    %% %%
    \bibitem{cppcoreguidelines-C21} \emph{If you define or =delete any copy, move, or destructor function, define or =delete them all}
    \newline
    Retreived from: \href{https://github.com/isocpp/CppCoreGuidelines/blob/master/CppCoreGuidelines.md#c21-if-you-define-or-delete-any-copy-move-or-destructor-function-define-or-delete-them-all}{Standard Cpp Foundation Github}

    %% %%
    \bibitem{isocpp-references} \emph{References}
    \newline
    Retrieved from: \href{https://isocpp.org/wiki/faq/references}{ISOC++ Wiki}

    %% %%
    \bibitem{isocpp-const-correctness} \emph{Const Correctness}
    \newline
    Retrieved from: \href{https://isocpp.org/wiki/faq/const-correctness}{ISOC++ Wiki}

    %% %%
    \bibitem{justswsolutions-lvalues_rvalues} (2016) - \emph{Core C++ - lvalues and rvalues}
    \newline
    Retrieved from: \href{https://www.justsoftwaresolutions.co.uk/cplusplus/core-c++-lvalues-and-rvalues.html}{Just Software Solutions}

    %% %%
    \bibitem{justswsolutions-rvalues_and_perfect_forwarding} (2008) - \emph{Rvalue References and Perfect Forwarding}
    \newline
    Retrieved from: \href{https://www.justsoftwaresolutions.co.uk/cplusplus/rvalue_references_and_perfect_forwarding.html}{Just Software Solutions}
    
    %% %%
    \bibitem{openstd-forwarding_problem} Peter Dimov, Howard E. Hinnant, Dave Abrahams - \emph{The Forwarding Problem: Arguments}
    \newline
    Retreived from: \href{https://www.open-std.org/jtc1/sc22/wg21/docs/papers/2002/n1385.htm}{open-std.org - N1385=02-0043}
    
    %% %%
    \bibitem{msa-adventure-in-perfect-forwarding}Meyers, Sutter, Alexandrescu (2011) - \emph{Session Announcement: Adventures in Perfect Forwarding}
    \newline
    Retrieved from: \href{https://cppandbeyond.com/2011/04/25/session-announcement-adventures-in-perfect-forwarding/}{C++ and Beyond}
    
    %% %%
    \bibitem{stackof-perfect-forwarding}\emph{What are the main purposes of std::forward and which problems does it solve?}
    \newline
    Retreived from: \href{https://stackoverflow.com/questions/3582001/what-are-the-main-purposes-of-stdforward-and-which-problems-does-it-solve}{StackOverflow Thread}

    %% %%
    \bibitem{stackof-move-semantic}\emph{What is the move semantic?}
    \newline
    Retreived from: \href{https://stackoverflow.com/questions/3106110/what-is-move-semantics}{StackOverflow Thread}

    %% %%
    \bibitem{slideshare-eywtk-about-move-semantic} Howard Hinnant (2014) - \emph{Everything you wanted to know about Move Semantics}
    \newline
    Retreived from: \href{https://www.slideshare.net/ripplelabs/howard-hinnant-accu2014}{Slideshare}
    
    %% %%
    \bibitem{cppcoreguidelines-F6}\emph{If your function must not throw declare it noexcept}
    \newline
    Retreived from: \href{https://github.com/isocpp/CppCoreGuidelines/blob/master/CppCoreGuidelines.md#f6-if-your-function-must-not-throw-declare-it-noexcept}{Cpp Core Guidelines}

    %% %%
    \bibitem{sof-auto-meaning}\emph{What is the meaning of \texttt{\textbf{auto}} keyword}
    \newline
    Retrieved from: \href{https://stackoverflow.com/questions/7576953/what-is-the-meaning-of-the-auto-keyword}{StackOverflow Thread}

    %% %%
    \bibitem{gotw-auto-variablesI} Herb Sutter - \emph{\texttt{\textbf{auto}} variables, Part 1} 
    \newline
    Retrieved from: \href{https://herbsutter.com/2013/06/07/gotw-92-solution-auto-variables-part-1/}{Guru of the weel - 92}

    %% %%
    \bibitem{gotw-auto-variablesII} Herb Sutter - \emph{\texttt{\textbf{auto}} variables, Part 2} 
    \newline
    Retrieved from: \href{https://herbsutter.com/2013/06/13/gotw-93-solution-auto-variables-part-2/}{Guru of the week - 93}

    %% %%
    \bibitem{gotw-almost-always-auto} Herb Sutter - \emph{AAA style (Almost Always \texttt{\textbf{auto}})}
    \newline
    Retrieved from: \href{https://herbsutter.com/2013/08/12/gotw-94-solution-aaa-style-almost-always-auto/}{Guru of the week - 94}

    %% %%
    \bibitem{gotw-interface-principle} - \emph{What's In a Class? - The Interface Principle}
    \newline
    Retrieved from: \href{http://www.gotw.ca/publications/mill02.htm}{Guru of the week - C++ Report, 10(3), March 1998}
    
    %% %%
    \bibitem{isocpp-non-type-template-arguments-auto}(2016)\emph{Declaring non-type template arguments with \texttt{\textbf{auto}}}
    \newline
    Retrieved from: \href{https://open-std.org/JTC1/SC22/WG21/docs/papers/2016/p0127r1.html}{Programming Language C++, Evolution Working Group - P0127R1}

    %% %%
    \bibitem{isocpp-non-type-template-param-auto}(2016)\emph{Declaring non-type template parameters with \texttt{\textbf{auto}}}
    \newline
    Retrieved from: \href{https://isocpp.org/files/papers/p0127r2.html}{Programming Language C++, Evolution Working Group - P0127R2}

    %% %%
    \bibitem{zemek-pros-and-cons-auto}\emph{Pros and Cons of Alternative Function Syntax in C++}
    \newline
    Retrieved from: \href{https://blog.petrzemek.net/2017/01/17/pros-and-cons-of-alternative-function-syntax-in-cpp/}{Petr Zemek's Blog of a Software Engineer}

    %% %%
    \bibitem{cppreference-adl}\emph{ADL - Argument-dependent lookup}
    \newline
    Retrieved from: \href{https://en.cppreference.com/w/cpp/language/adl}{Cppreference}

    %% %%
    \bibitem{cppreference-lambda}\emph{Lambda Expressions}
    \newline
    Retrieved from: \href{https://en.cppreference.com/w/cpp/language/lambda}{Cppreference}
    
    %% %%
    \bibitem{openstd-wording-morph-lambda}Larvi, Freeman, Crowl (2008) - \emph{Lambda Expressions and Closures: Wording for Monomorphic Lambdas (Revision 4)}
    \newline
    Retrieved from: \href{https://www.open-std.org/jtc1/sc22/wg21/docs/papers/2008/n2550.pdf}{open-std.org}

    %% %%
    \bibitem{openstd-wording-C++0x}Vandevoorde (2009) - \emph{New wording for C++0x Lambdas}
    \newline
    Retrieved from: \href{https://www.open-std.org/jtc1/sc22/wg21/docs/papers/2009/n2859.pdf}{open-std.org}      
    
    %% %%
    \bibitem{wikipedia-closure}\emph{Closure (Computer Programming)}
    \newline
    Retrieved from: \href{https://en.wikipedia.org/wiki/Closure_(computer_programming)}{WikiPedia}

    %% %%
    \bibitem{cppreference-abbreviated-functemplate}\emph{Abbreviated function template}
    \newline
    Retrieved from: \href{https://en.cppreference.com/w/cpp/language/function_template#Abbreviated_function_template}{Cppreference}

    %% %%
    \bibitem{cppreference-decltype}\texttt{decltype} \emph{specifier}
    \newline
    Retrieved from: \href{https://en.cppreference.com/w/cpp/language/decltype}{Cppreference}

    %% %%
    \bibitem{openstd-initializer-lists-rev3}Stroustrup, Dos Reis (2007) - \emph{Initializer list (Rev. 3)}
    \newline
    Retrieved from: \href{https://www.open-std.org/jtc1/sc22/wg21/docs/papers/2007/n2215.pdf}{open-std.org}
    
    %% %%
    \bibitem{openstd-initializer-lists-alt-mech-ratio}Merrill, Vandevoorde (2008) - \emph{Initializer Lists - Alternative Mechanism and Rationale (v.2)}
    \newline
    Retrieved from: \href{https://www.open-std.org/jtc1/sc22/wg21/docs/papers/2008/n2640.pdf}{Initializer Lists — Alternative Mechanism and Rationale (v. 2}
    
    %% %%
    \bibitem{openstd-range-for-loop-rev2}Thorsten Ottosen (2007) - \emph{Wording for range-based for-loop (revision 2)}
    \newline
    Retrieved from: \href{https://www.open-std.org/jtc1/sc22/wg21/docs/papers/2007/n2243.html}{open-std.org}    
    
    %% %%
    \bibitem{openstd-range-based-for-adl}\emph{Range-based \texttt{for} statements and ADL}
    \newline
    Retrieved from: \href{https://www.open-std.org/jtc1/sc22/wg21/docs/papers/2011/n3257.pdf}{open-std.org}    
    
    %% %%
    \bibitem{openstd-delete-op-default-noexcept}Svoboda (2010) - \emph{Delete operators default to noexcept}
    \newline
    Retrieved from: \href{https://www.open-std.org/jtc1/sc22/wg21/docs/papers/2010/n3205.htm}{open-std.org}
    
    %% %%
    \bibitem{openstd-noexcept-for-dctors}Mauer (2010) - \emph{Deducing "noexcept" for destructors}
    \newline
    Retrieved from: \href{https://www.open-std.org/jtc1/sc22/wg21/docs/papers/2010/n3204.htm}{open-std.org}
    
    %% %%
    \bibitem{openstd-move-ctor-and-throw}Abrahams, Sharoni, Gregor (2010) - \emph{Allowing Move Constructors to Throw (Rev. 1)}
    \newline
    Retrieved from: \href{https://www.open-std.org/jtc1/sc22/wg21/docs/papers/2010/n3050.html}{open-std.org}
    
    %~ %% %%
    \bibitem{stroustrup-C++-faq}\emph{Bjarne Stroustrup's C++ Style and Technique FAQ}
    \newline
    Retrieved from: \href{https://www.stroustrup.com/bs_faq2.html}{stroutrup.com}    
    
    %~ %% %%
    %~ \bibitem{}\emph{}
    %~ \newline
    %~ Retrieved from: \href{}{}    
    % C++ Chapter biblio ends
    %%%
%%%
